\begin{theorem}\label{thm:p143-mp-section}
Let \(X\) be a connected compact K\"ahler manifold, let
\(\alpha=\{\omega\}\) be a K\"ahler class, and assume the all-point
upper-semicontinuity conclusion of
\Cref{thm:p142-raw-transform-usc,cor:p142-pseudoeffective-raw-usc}.
Then the raw ideal-valuative transform has a unique section
\[
 \mathcal S_\alpha:\PSH_{\mathrm{MP}}(\alpha)
        \longrightarrow\PSH^{\NA}_{\mathrm{DXZ}}(X,\omega),
 \qquad
 \mathfrak T_{\mathcal S_\alpha(\phi)}=\phi
 \quad\text{on }X^{\bethan}.
\]
The section preserves and reflects pointwise order.

More precisely, let \((\phi_i)_{i\in I}\) be any decreasing net in
\(\PL(X^{\bethan})\cap\PSH_{\mathrm{MP}}(\alpha)\) with pointwise limit
\(\phi\not\equiv-\infty\).  For every \(i\), let
\(\boldsymbol\Gamma^i\in
\mathcal E^{1,\NA}_{\mathrm{DXZ}}(X,\omega)\) be the object corresponding
to \(\phi_i\) under \eqref{eq:pseff-finite-energy-bijection}.  Then the
\(\boldsymbol\Gamma^i\) form a decreasing net,
\[
                 \inf_i\boldsymbol\Gamma^i_{\max}>-\infty,
\]
and
\[
                 \mathcal S_\alpha(\phi)
                 =\inf_i\boldsymbol\Gamma^i.
\]
The infimum and its raw transform are independent of the chosen model net.
Thus every MP \(\alpha\)-psh function, including one not known to have
finite energy, belongs to the image of the raw transform.
\end{theorem}

\begin{remark}\label{rem:p143-mp-section-scope}
The theorem does not prove that the raw transform of every DXZ object is
MP-psh.  Equivalently, it identifies a canonical copy of the whole MP cone
inside the DXZ space but does not show that this copy exhausts the DXZ
space.  The missing direction is a cofinal positivity statement for an
arbitrary DXZ object, not another divisorial-density argument.
\end{remark}
